\documentclass[11pt]{article}

\usepackage[a4paper]{geometry}
\geometry{left=2.0cm,right=2.0cm,top=2.5cm,bottom=2.5cm}

% Ensure only one \begin{document} and all environments are closed properly

\usepackage{ctex} % 支持中文的LaTeX宏包
\usepackage{amsmath,amsfonts,graphicx,subfigure,amssymb,bm,amsthm,mathrsfs,mathtools,breqn} % 数学公式和符号的宏包集合
\usepackage{algorithm,algorithmicx} % 算法和伪代码
\usepackage[noend]{algpseudocode} % 算法和伪代码
\usepackage{fancyhdr} % 自定义页眉页脚
\usepackage[framemethod=TikZ]{mdframed} % 创建带边框的框架
\usepackage{fontspec} % 字体设置
\usepackage{adjustbox} % 调整盒子大小
\usepackage{fontsize} % 设置字体大小
\usepackage{tikz,xcolor} % 绘制图形和使用颜色
\usepackage{multicol} % 多栏排版
\usepackage{multirow} % 表格中合并单元格
\usepackage{pdfpages} % 插入PDF文件
\usepackage{listings} % 在文档中插入源代码
\usepackage{wrapfig} % 文字绕排图片
\usepackage{bigstrut,multirow,rotating} % 支持在表格中使用特殊命令
\usepackage{booktabs} % 创建美观的表格
\usepackage{circuitikz} % 绘制电路图
\usepackage{zhnumber} % 中文序号（用于标题）
\usepackage{tabularx} % 表格折行
\usepackage{tikz}
\usetikzlibrary{positioning} % 加载 positioning 库

\definecolor{dkgreen}{rgb}{0,0.6,0}
\definecolor{gray}{rgb}{0.5,0.5,0.5}
\definecolor{mauve}{rgb}{0.58,0,0.82}
\lstset{
  frame=tb,
  aboveskip=3mm,
  belowskip=3mm,
  showstringspaces=false,
  columns=flexible,
  framerule=1pt,
  rulecolor=\color{gray!35},
  backgroundcolor=\color{gray!5},
  basicstyle={\small\ttfamily},
  numbers=none,
  numberstyle=\tiny\color{gray},
  keywordstyle=\color{blue},
  commentstyle=\color{dkgreen},
  stringstyle=\color{mauve},
  breaklines=true,
  breakatwhitespace=true,
  tabsize=3,
}

% 超链接支持 (一定要在大多数包之后, cleveref 之前)
% colorlinks=true 使链接文字着色而不是加方框; hidelinks 可去掉颜色
\usepackage[colorlinks=true,linkcolor=blue,citecolor=blue,urlcolor=magenta]{hyperref}

% 轻松引用, 可以用\cref{}指令直接引用, 自动加前缀. 需要在 hyperref 之后加载
% 例: 图片label为fig:1  => \cref{fig:1} -> Figure.1; \ref{fig:1} -> 1 (纯数字)
\usepackage[capitalize]{cleveref}
% \crefname{section}{Sec.}{Secs.}
\Crefname{section}{Section}{Sections}
\Crefname{table}{Table}{Tables}
\crefname{table}{Table.}{Tabs.}

% \setmainfont{Palatino Linotype.ttf}
% \setCJKmainfont{SimHei.ttf}
% \setCJKsansfont{Songti.ttf}
% \setCJKmonofont{SimSun.ttf}
\punctstyle{kaiming}
% 偏好的几个字体, 可以根据需要自行加入字体ttf文件并调用

\renewcommand{\emph}[1]{\begin{kaishu}#1\end{kaishu}}

% 对 section 等环境的序号使用中文
\renewcommand \thesection{\zhnum{section}、}
\renewcommand \thesubsection{\arabic{section}.\arabic{subsection}}

\tikzset{
    link/.style={draw, thick, -stealth} % 定义 link 样式，包含绘制、粗线条和箭头
}


%%%%%%%%%%%%%%%%%%%%%%%%%%%
%改这里可以修改实验报告表头的信息
\newcommand{\name}{寇逸欣}
\newcommand{\studentNum}{2023K8009922004}
\newcommand{\major}{计算机科学与技术}
\newcommand{\labNum}{14}
\newcommand{\labName}{网络传输机制实验-可靠传输}
%%%%%%%%%%%%%%%%%%%%%%%%%%%

\begin{document}

\input{tex_file/head.tex}

\section{实验内容}
\begin{enumerate}
    \item 执行 \texttt{create\_randfile.sh}，生成待传输数据文件 \texttt{client-input.dat}
    \item 运行给定网络拓扑（\texttt{tcp\_topo\_loss.py}）
    \item 在节点 \texttt{h1} 上执行 TCP 程序
    \begin{enumerate}
        \item 执行脚本（\texttt{disable\_offloading.sh}，\texttt{disable\_tcp\_rst.sh}），禁止协议栈的相应功能
        \item 在 \texttt{h1} 上运行 TCP 协议栈的服务器模式（\texttt{./tcp\_stack server 10001}）
    \end{enumerate}
    \item 在节点 \texttt{h2} 上执行 TCP 程序
    \begin{enumerate}
        \item 执行脚本（\texttt{disable\_offloading.sh}，\texttt{disable\_tcp\_rst.sh}），禁止协议栈的相应功能
        \item 在 \texttt{h2} 上运行 TCP 协议栈的客户端模式（\texttt{./tcp\_stack client 10.0.0.1 10001}）
        \begin{enumerate}
            \item 客户端发送文件 \texttt{client-input.dat} 给服务器，服务器将收到的数据存储到文件 \texttt{server-output.dat}
        \end{enumerate}
    \end{enumerate}
    \item 使用 \texttt{md5sum} 比较两个文件是否完全相同
    \item 使用 \texttt{tcp\_stack.py} 替换两端任意一方，对端都能正确处理数据收发
\end{enumerate}

\begin{figure}[H]
    \centering
    \includegraphics[width=0.5\textwidth]{fig/state_machine.png}
    \caption{TCP 状态机图}
    \label{fig:state_machine}
\end{figure}

\section{实验步骤}

\subsection{数据结构设计}
为了支持发送队列和乱序接收队列的管理，定义了新的数据结构 \texttt{struct data\_packet}，用于存储数据包的相关信息。同时在 \texttt{struct tcp\_sock} 中增加了发送队列 \texttt{send\_buf} 和乱序接收队列 \texttt{rcv\_ofo\_buf} 及其对应的互斥锁。

\begin{lstlisting}[language=C, caption=数据结构定义]
struct data_packet {
    struct list_head list;
    u8 flags;
    u8 times;       // 重传次数
    u32 seq;
    u32 len;
    u32 seq_end;
    char *packet;   // 数据包内容
};

struct tcp_sock {
    // ... existing code ...
    // used to pend unacked packets
    struct list_head send_buf;
    // used to pend out-of-order packets
    struct list_head rcv_ofo_buf;

    pthread_mutex_t send_buf_lock;
    pthread_mutex_t rcv_buf_lock;
    // ... existing code ...
};
\end{lstlisting}

此外，为了支持微秒级的重传定时器，修改了 \texttt{struct tcp\_timer} 中的 \texttt{timeout} 字段为 \texttt{long long} 类型。

\subsection{发送队列维护}
\subsubsection{加入发送队列}
实现了 \texttt{new\_data\_block} 函数用于构建 \texttt{data\_packet} 结构体。在 \texttt{tcp\_sock\_connect} (SYN)、\texttt{tcp\_sock\_close} (FIN) 和 \texttt{tcp\_sock\_write} (Data) 等发送数据的函数中，调用该函数将数据包加入 \texttt{send\_buf} 队列，并启动重传定时器。

\begin{lstlisting}[language=C]
// 示例：在 tcp_sock_write 中加入发送队列
struct data_packet *dp = new_data_block(TCP_PSH|TCP_ACK, tsk->snd_nxt, data_len, buf + len_send);
pthread_mutex_lock(&tsk->send_buf_lock);
list_add_tail(&dp->list, &tsk->send_buf);
pthread_mutex_unlock(&tsk->send_buf_lock);

// 发送数据包
tcp_send_packet(tsk, packet, packet_len);

// 启动重传定时器
if(!tsk->retrans_timer.enable)
    tcp_set_retrans_timer(tsk);
\end{lstlisting}

\subsubsection{清除已确认数据}
实现了 \texttt{tcp\_free\_send\_buf} 函数。当收到 ACK 时，在 \texttt{tcp\_process} 中调用该函数，遍历 \texttt{send\_buf}，删除 \texttt{seq\_end} 小于等于 \texttt{ack} 的数据包，并更新定时器状态。

\begin{lstlisting}[language=C]
void tcp_free_send_buf(struct tcp_sock *tsk, struct tcp_cb *cb)
{
    struct data_packet *tmp, *q;
    pthread_mutex_lock(&tsk->send_buf_lock);
    list_for_each_entry_safe(tmp, q, &tsk->send_buf, list)
    {
        if(less_or_equal_32b(tmp->seq_end, cb->ack))
        {
            list_delete_entry(&tmp->list);
            if(tmp->packet) free(tmp->packet);
            free(tmp);
            // 重置定时器逻辑...
        }
    }
    pthread_mutex_unlock(&tsk->send_buf_lock);
}
\end{lstlisting}

\subsection{乱序接收处理}
\subsubsection{缓存乱序包}
实现了 \texttt{tcp\_rcv\_ofo\_pkt} 函数。当接收到的数据包序列号不等于 \texttt{rcv\_nxt} 时，调用该函数将数据包按序列号顺序插入 \texttt{rcv\_ofo\_buf} 队列，并回复 ACK（携带期望的 \texttt{rcv\_nxt}）。

\begin{lstlisting}[language=C]
void tcp_rcv_ofo_pkt(struct tcp_sock *tsk, struct tcp_cb *cb)
{
    struct data_packet *dp = new_data_block(cb->flags, cb->seq, cb->pl_len, cb->payload);
    // ... 遍历队列找到合适位置插入 ...
    list_add_tail(&dp->list, pos);
}
\end{lstlisting}

\subsubsection{处理接收队列}
在 \texttt{tcp\_process} 中，当收到符合 \texttt{rcv\_nxt} 的数据包后，将其写入接收缓冲区，随后检查 \texttt{rcv\_ofo\_buf}。如果队列头部的数据包序列号等于更新后的 \texttt{rcv\_nxt}，则将其取出并写入接收缓冲区，重复此过程直到队列为空或序列号不连续。

\begin{lstlisting}[language=C]
// 检查乱序队列
list_for_each_entry_safe(tmp, q, &tsk->rcv_ofo_buf, list) {
    if(tmp->seq == tsk->rcv_nxt) {
        // ... 写入接收缓冲区 ...
        tsk->rcv_nxt = tmp->seq_end;
        // ... 处理 FIN 标志 ...
        list_delete_entry(&tmp->list);
        free(tmp);
    } else {
        break;
    }
}
\end{lstlisting}

\subsection{超时重传机制}
\subsubsection{定时器管理}
实现了 \texttt{tcp\_set\_retrans\_timer} 和 \texttt{tcp\_unset\_retrans\_timer} 函数用于开启和关闭重传定时器。定时器类型设为 1。

\subsubsection{超时处理与指数避退}
修改了 \texttt{tcp\_scan\_timer\_list} 函数。当检测到重传定时器超时时：
1. 检查重传次数，若超过 3 次，则发送 RST 并关闭连接。
2. 若未超过次数，则重传队首数据包，并将超时时间翻倍（指数避退）。
3. 调用 \texttt{tcp\_send\_retrans\_packet} 发送数据包（不增加 \texttt{snd\_nxt}）。

\begin{lstlisting}[language=C]
// 指数避退计算
long long duration = TCP_RETRANS_INTERVAL_INITIAL * (1 << dp->times);
if (cur_time - entry->timeout > duration) {
    if (dp->times >= 3) {
        // ... 发送 RST 并关闭连接 ...
    } else {
        dp->times++;
        entry->timeout = cur_time;
        tcp_send_retrans_packet(tsk, dp);
    }
}
\end{lstlisting}

\section{实验结果}
先调用 \texttt{create\_randfile.sh} 生成待传输数据文件 \texttt{client-input.dat}，然后运行给定的网络拓扑脚本 \texttt{tcp\_topo\_loss.py}。在节点 \texttt{h1} 上执行 TCP 协议栈的服务器模式，并在节点 \texttt{h2} 上执行客户端模式。客户端发送文件后，服务器将收到的数据存储到文件 \texttt{server-output.dat}。
\begin{figure}[H]
    \centering
    \includegraphics[width=0.8\textwidth]{fig/h1_log.png}
    \caption{h1 log 输出}
    \label{fig:h1_log}
\end{figure}

\begin{figure}[H]
    \centering
    \includegraphics[width=0.8\textwidth]{fig/h2_log.png}
    \caption{h2 log 输出}
    \label{fig:h2_log}
\end{figure}

使用 \texttt{md5sum} 比较两个文件，结果显示两者完全相同，验证了可靠传输的正确性。
\begin{figure}[H]
    \centering
    \includegraphics[width=0.5\textwidth]{fig/md5sum.png}
    \caption{md5sum 比较结果}
    \label{fig:md5sum}
\end{figure}

\section{实验总结}
本次实验代码量并不大，难点在于如何构建发送队列和乱序接收队列，以及如何处理超时重传机制（包括状态机的修改，缓存包和释放的时机）。实验结果表明，修改后的 TCP 协议栈能够正确处理数据传输，并且在丢包和乱序情况下仍能保证数据的完整性和顺序。

在这次实验中遇到了一个重复释放的问题，这是由于在处理 TCP_LISTEN 状态收到 SYN 包时，使用了 \texttt{memcpy} 将父 socket（\texttt{tsk}）的内容复制到了新分配的子 socket（\texttt{child\_tsk}）中引起的。\texttt{alloc\_tcp\_sock} 为 \texttt{child\_tsk} 分配的资源被 \texttt{memcpy} 覆盖，导致 \texttt{child\_tsk} 和 \texttt{tsk} 共享同一个 \texttt{rcv\_buf}，从而引发了双重释放错误。解决方案是避免使用 \texttt{memcpy}，手动初始化子 socket，仅复制必要字段，确保父子 socket 的资源独立。
\end{document}